\documentclass[12pt, a4paper]{article}

% --- Pacotes ---
\usepackage[utf8]{inputenc}
\usepackage[T1]{fontenc}
\usepackage[brazil]{babel}
\usepackage{geometry}
\usepackage{setspace}
\usepackage{graphicx}
\usepackage{amsmath}
\usepackage{amssymb}
\usepackage{hyperref}
\usepackage{booktabs}
\usepackage{array}
\usepackage{caption}
\usepackage{subcaption}
\usepackage{cite}
\usepackage{indentfirst}
\usepackage{xcolor}
\usepackage{fancyhdr}
\usepackage{titlesec}
\usepackage{abstract}

\geometry{
  left=3cm,
  right=2cm,
  top=3cm,
  bottom=2cm
}

\onehalfspacing

% --- Cabeçalho e rodapé ---
\pagestyle{fancy}
\fancyhf{}
\rhead{\small Previsão de Vida Útil de Baterias}
\lhead{\small Resenha de Artigo Científico}
\cfoot{\thepage}

% --- Formatação de títulos ---
\titleformat{\section}{\large\bfseries}{\thesection}{1em}{}
\titleformat{\subsection}{\normalsize\bfseries}{\thesubsection}{1em}{}

\hypersetup{
  colorlinks=true,
  linkcolor=blue!70!black,
  citecolor=blue!70!black,
  urlcolor=blue!70!black
}

% ================================================================
\begin{document}

% --- Capa ---
\begin{titlepage}
  \centering
  \vspace*{2cm}

  {\Large \textbf{Universidade --- Programa de Pós-Graduação}}\\[0.5cm]
  {\large Doutorado em Engenharia / Área de Concentração}\\

  \vspace{3cm}

  {\LARGE \textbf{Resenha Crítica de Artigo Científico}}\\[1cm]

  {\Large \textit{Data-Driven Prediction of Battery Cycle Life\\
  Before Capacity Degradation}}\\[0.5cm]

  {\large Severson, Attia \textit{et al.} --- \textit{Nature Energy}, 2019}

  \vspace{3cm}

  \begin{flushright}
    \textbf{Autor(a):} Pedro Lievra\\[0.3cm]
    \textbf{E-mail:} \href{mailto:pedrolievra@gmail.com}{pedrolievra@gmail.com}\\[0.3cm]
    \textbf{Orientador(a):} [Nome do Orientador]\\[0.3cm]
    \textbf{Data:} Junho de 2026
  \end{flushright}

  \vfill
  {\small Documento gerado para fins acadêmicos de doutorado.}
\end{titlepage}

% --- Resumo ---
\begin{abstract}
  \noindent
  Esta resenha analisa o artigo \textit{``Data-Driven Prediction of Battery
  Cycle Life Before Capacity Degradation''} publicado na revista
  \textit{Nature Energy} em 2019 por Severson, Attia e colaboradores. O
  trabalho propõe uma abordagem de aprendizado de máquina capaz de prever
  quantitativamente a vida útil de baterias de íon-lítio a partir dos
  ciclos iniciais de operação --- muito antes de qualquer degradação visível
  de capacidade. São discutidos o contexto do problema, a construção do
  conjunto de dados, as limitações das features convencionais baseadas em
  capacidade de descarga e a motivação para o uso de curvas de tensão
  diferencial como variável preditora. A análise cobre o artigo até o
  tópico \textit{Machine-Learning Approach}, incluindo a descrição das
  figuras relevantes.

  \medskip
  \noindent\textbf{Palavras-chave:} baterias de íon-lítio, previsão de vida
  útil, aprendizado de máquina, degradação eletroquímica, fast-charging, LFP.
\end{abstract}

\newpage
\tableofcontents
\newpage

% ================================================================
\section{Introdução e Contextualização}

Baterias de íon-lítio são componentes críticos em uma vasta gama de
aplicações modernas: veículos elétricos, dispositivos portáteis e sistemas
de armazenamento de energia em larga escala. Nessas aplicações, a vida útil
da bateria --- tipicamente definida como o número de ciclos de carga e
descarga até que a capacidade caia para 80\% do valor nominal --- é um
parâmetro de projeto fundamental \cite{severson2019}.

A principal dificuldade em prever esse parâmetro reside na natureza da
degradação: ela é eletroquimicamente não-linear, fortemente dependente das
condições operacionais (temperatura, protocolo de carga, profundidade de
descarga) e praticamente \emph{imperceptível} nos ciclos iniciais. Em
termos práticos, isso significa que o feedback sobre o desempenho real de
uma bateria só se materializa após meses ou anos de uso --- um obstáculo
severo para o desenvolvimento acelerado de novos materiais e protocolos de
carregamento.

Trabalhos anteriores de aprendizado de máquina aplicados a baterias
esbarravam nessa limitação. De modo geral, previsões confiáveis exigiam
dados correspondentes a pelo menos 25\% do caminho até a falha, ou então
dependiam de medições especializadas e de difícil obtenção em campo
\cite{severson2019}. Harris \textit{et al.}, por exemplo, relataram
correlação muito fraca ($\rho = 0{,}1$) entre a capacidade no ciclo 80 e a
capacidade no ciclo 500, evidenciando que as features mais óbvias não
carregam informação preditiva suficiente nos estágios iniciais.

O artigo de Severson, Attia \textit{et al.} \cite{severson2019} apresenta
uma resposta a esse desafio: demonstra que é possível prever a vida útil
quantitativamente com \textbf{9,1\% de erro de teste} usando dados dos
primeiros 100 ciclos, período em que a variação mediana de capacidade das
células é de apenas $+0{,}2\%$ em relação ao valor inicial. Adicionalmente,
uma tarefa de classificação binária (vida curta vs.\ vida longa) atinge
\textbf{4,9\% de erro} utilizando apenas os primeiros 5 ciclos. A chave
para esses resultados é o uso de curvas de voltagem diferenciais ao invés
da capacidade de descarga simples.

% ================================================================
\section{O Conjunto de Dados}

\subsection{Células e Protocolo Experimental}

Para viabilizar o estudo, os autores conduziram um experimento controlado em
larga escala. Foram cicladas \textbf{124 células comerciais} do modelo
A123~APR18650M1A (química LFP/grafite, capacidade nominal 1,1~Ah), todas
nominalmente idênticas. As células foram submetidas a \textbf{72 protocolos
de carregamento rápido} (\textit{fast-charging}) distintos, com taxas de
corrente variando de 3,6C a 6C e com descargas todas idênticas a 4C até
2,0~V. Todas as células foram mantidas em câmara climatizada a 30~°C.

\subsection{Características do Dataset}

O resultado é o maior conjunto de dados público de baterias comerciais
nominalmente idênticas cicladas sob condições controladas até a data de
publicação. As principais características estão sumarizadas na
Tabela~\ref{tab:dataset}.

\begin{table}[ht]
  \centering
  \caption{Características do dataset de Severson \textit{et al.} (2019).}
  \label{tab:dataset}
  \begin{tabular}{lc}
    \toprule
    \textbf{Característica} & \textbf{Valor} \\
    \midrule
    Total de células             & 124 \\
    Lotes (\textit{batches})     & 3 (46, 48 e 45 células) \\
    Total de ciclos acumulados   & $\approx$ 96.700 \\
    Vida útil mínima             & $\sim$150 ciclos \\
    Vida útil máxima             & $\sim$2.300 ciclos \\
    Vida útil média              & 806 ciclos \\
    Desvio-padrão da vida útil   & 377 ciclos \\
    Temperatura de operação      & 30~°C \\
    Critério de fim de vida (EOL)& 80\% da capacidade nominal \\
    \bottomrule
  \end{tabular}
\end{table}

A ampla variabilidade nas vidas úteis (de $\sim$150 a $\sim$2.300 ciclos,
fator $\approx$15) sob condições controladas é notável e sublinha a
sensibilidade do processo de degradação às condições de carga, mesmo entre
células nominalmente idênticas. Esse espectro amplo é justamente o que torna
o problema interessante e o dataset valioso para o treinamento de modelos
preditivos.

\subsection{Estrutura dos Dados}

Para cada célula, o dataset registra três categorias de dados:

\begin{itemize}
  \item \textbf{Descritores:} política de carregamento, vida útil medida,
  código de barras e canal do equipamento.

  \item \textbf{Dados de resumo por ciclo (\textit{summary}):} capacidade
  de descarga ($Q_D$), capacidade de carga ($Q_C$), resistência interna
  (IR), temperaturas (máxima, média e mínima) e tempo de carga até 80\%
  do estado de carga (SOC).

  \item \textbf{Dados dentro do ciclo (\textit{cycle data}):} séries
  temporais de corrente ($I$), voltagem ($V$), temperatura ($T$),
  capacidades acumuladas de carga e descarga, e a curva diferencial
  $dQ/dV$ durante a descarga.
\end{itemize}

% ================================================================
\section{Limitações das Features Baseadas em Capacidade}

\subsection{Motivação para uma Abordagem Alternativa}

A seção mais ilustrativa da primeira parte do artigo investiga por que a
medida mais natural --- a capacidade de descarga --- falha como preditor de
vida útil nos ciclos iniciais.

Os autores calcularam o coeficiente de correlação de Pearson ($\rho$) entre
o logaritmo da vida útil e três features derivadas da curva de capacidade:

\begin{enumerate}
  \item \textbf{Capacidade de descarga no ciclo 2:} $\rho = -0{,}06$
  (essencialmente sem correlação).
  \item \textbf{Capacidade de descarga no ciclo 100:} $\rho = 0{,}27$
  (correlação fraca).
  \item \textbf{Inclinação da curva de capacidade nos ciclos 95--100:}
  $\rho = 0{,}47$ (melhor resultado, mas ainda modesto).
\end{enumerate}

Nenhuma dessas features oferece poder preditivo suficiente. Este resultado
é central para o argumento do artigo: a informação relevante sobre o
mecanismo de degradação não está exposta na curva de capacidade durante os
ciclos iniciais.

\subsection{O Fenômeno de Aumento de Capacidade}

Um resultado contraintuitivo reportado pelos autores é que, em
\textbf{81\% das células}, a capacidade de descarga no ciclo 100 é
\emph{maior} do que no ciclo 2. Esse comportamento é atribuído ao fenômeno
de redistribuição de carga: parte da capacidade nominalmente armazenada na
região do eletrodo negativo que se estende além dos limites do positivo
(\textit{overhang}) gradualmente se equaliza durante os ciclos iniciais,
elevando aparentemente a capacidade útil.

Esse efeito mascara a degradação real que está ocorrendo em escala
microscópica (formação e crescimento da camada de interface sólido-eletrolítica,
SEI; perda de lítio ciclável), tornando a capacidade de descarga um proxy
enganoso nos primeiros ciclos.

% ================================================================
\section{Descrição das Figuras (até a abordagem de ML)}

\subsection{Figura 1 --- Baixo Poder Preditivo das Features de Capacidade}

A Figura~1 do artigo sintetiza visualmente os argumentos da Seção~3 desta
resenha. Ela é composta por seis painéis (a--f), todos com a codificação de
cor consistente utilizada no restante do artigo: cores mais escuras
(azul/roxo) representam células de vida longa; cores mais claras
(amarelo/vermelho) representam células de vida curta.

\paragraph{Painel (a) --- Capacidade de descarga vs.\ número de ciclo
(primeiros 1.000 ciclos).}
Exibe as curvas de capacidade de descarga para todas as 124 células ao
longo dos primeiros 1.000 ciclos. O comportamento clássico de degradação de
baterias LFP é evidente: a capacidade permanece relativamente estável por
um longo período e depois cai abruptamente próximo ao fim de vida
(\textit{``knee point''}). O aspecto mais importante é o entrelaçamento
das curvas: células que iniciam com maior capacidade não necessariamente
vivem mais, e vice-versa. Isso já antecipa visualmente a baixa correlação
quantificada nos painéis seguintes.

\paragraph{Painel (b) --- Zoom dos primeiros 100 ciclos.}
Ampliação do painel (a) restrita ao intervalo de 0 a 100 ciclos. As curvas
estão completamente misturadas numa faixa estreita, sem qualquer separação
visual entre células de vida curta e longa. É o argumento visual mais direto
para a impossibilidade de prever a vida útil usando apenas a trajetória de
capacidade nesse intervalo.

\paragraph{Painel (c) --- Histograma da razão de capacidade (ciclo 100 /
ciclo 2).}
A linha pontilhada vertical em 1,00 divide as células entre aquelas que
ganharam e aquelas que perderam capacidade até o ciclo 100. A grande maioria
das células está à direita de 1,00 (ganho de capacidade), ilustrando o
fenômeno de aumento de capacidade discutido anteriormente. Uma célula
com razão $\approx 0{,}90$ foi excluída da visualização para melhor
resolução do histograma.

\paragraph{Painéis (d), (e) e (f) --- Gráficos de dispersão: vida útil
vs.\ features de capacidade.}
Esses três painéis apresentam, em escala logarítmica no eixo $y$, a
correlação entre a vida útil e cada uma das três features:

\begin{itemize}
  \item[(d)] Capacidade no ciclo 2: nuvem de pontos sem tendência ($\rho = -0{,}06$).
  \item[(e)] Capacidade no ciclo 100: tendência ligeiramente positiva ($\rho = 0{,}27$).
  \item[(f)] Inclinação da capacidade entre os ciclos 95--100
  (mAh/ciclo): correlação moderada ($\rho = 0{,}47$), a melhor
  entre as três, mas ainda insuficiente para previsão confiável.
\end{itemize}

A progressão $(\text{d}) \to (\text{e}) \to (\text{f})$ demonstra que,
mesmo tentando extrair features mais sofisticadas da curva de capacidade, o
poder preditivo permanece baixo. Essa limitação fundamenta a necessidade de
uma feature alternativa: a curva diferencial de capacidade $\Delta Q(V)$,
apresentada na seção de aprendizado de máquina do artigo.

% ================================================================
\section{Considerações sobre a Metodologia Experimental}

\subsection{Pontos Fortes do Delineamento}

O delineamento experimental adotado apresenta características que elevam
significativamente a relevância científica dos resultados:

\begin{itemize}
  \item \textbf{Controle de variáveis confundidoras:} ao utilizar células
  nominalmente idênticas, temperatura constante e protocolo de descarga
  uniforme, os autores isolam o efeito do protocolo de carga como fonte
  de variabilidade nas vidas úteis.

  \item \textbf{Escala do dataset:} 96.700 ciclos acumulados constitui um
  dataset de magnitude incomum em pesquisa de baterias, onde experimentos
  são caros e demorados.

  \item \textbf{Reprodutibilidade:} os dados foram disponibilizados
  publicamente em \url{https://data.matr.io/1/}, permitindo validação
  independente e reutilização por outros grupos.

  \item \textbf{Separação clara de conjuntos de treino/teste:} os autores
  definiram o split treino/teste de forma explícita e replicável (células
  ímpares dos Lotes 1 e 2 para teste; todas as células do Lote 3 como
  conjunto de teste secundário).
\end{itemize}

\subsection{Limitações e Perspectivas}

Apesar dos resultados expressivos, algumas limitações merecem menção:

\begin{itemize}
  \item O dataset utiliza apenas uma química de eletrodo (LFP/grafite) e
  uma geometria de célula (cilíndrica 18650). A generalização para NMC,
  NCA ou células de formato \textit{pouch} requer validação adicional.

  \item O código de modelagem preditiva está sob licença acadêmica
  restrita (contato: braatz@mit.edu), o que limita a reprodutibilidade
  plena dos resultados de aprendizado de máquina.

  \item O ambiente controlado (30~°C, câmara climática) pode não capturar
  a variabilidade térmica de aplicações reais, onde flutuações de
  temperatura são relevantes para a degradação.
\end{itemize}

% ================================================================
\section{Conclusão Parcial}

Até o ponto ``Machine-Learning Approach'', o artigo de Severson \textit{et
al.} constrói um argumento rigoroso e bem fundamentado: a capacidade de
descarga, a métrica mais intuitiva para acompanhar o estado de saúde de uma
bateria, é um preditor \emph{fraco} durante os ciclos iniciais. O
aumento de capacidade observado em 81\% das células nos primeiros 100 ciclos
é uma evidência concreta de que fenômenos de redistribuição interna de carga
mascaram a degradação real.

O dataset construído para o estudo --- 124 células, $\sim$96.700 ciclos,
disponível publicamente --- representa por si só uma contribuição
significativa para a comunidade científica. A análise das correlações entre
features de capacidade e vida útil ($\rho$ máximo de 0,47) demonstra
quantitativamente a necessidade de uma abordagem diferente, pavimentando o
caminho para a feature $\Delta Q(V)$ e o modelo de aprendizado de máquina
que são discutidos na sequência do artigo.

Os resultados finais (9,1\% de erro para previsão quantitativa com 100
ciclos; 4,9\% para classificação binária com apenas 5 ciclos) representam
um avanço substancial em relação ao estado da arte e têm implicações diretas
para o desenvolvimento de protocolos de fast-charging, gerenciamento de
baterias e avaliação de degradação em tempo real.

% ================================================================
\section*{Referências}
\addcontentsline{toc}{section}{Referências}

\begin{thebibliography}{9}

\bibitem{severson2019}
Severson, K.~A., Attia, P.~M., Jin, N., Perkins, N., Jiang, B.,
Yang, Z., Chen, M.~H., Aykol, M., Herring, P.~K., Fraggedakis, D.,
Bazant, M.~Z., Harris, S.~J., Chueh, W.~C., \& Braatz, R.~D. (2019).
\textit{Data-driven prediction of battery cycle life before capacity
degradation.}
\textit{Nature Energy}, 4(5), 383--391.
\url{https://doi.org/10.1038/s41560-019-0356-8}

\bibitem{harris2017}
Harris, S.~J., Harris, D.~J., \& Li, C. (2017).
\textit{Failure statistics for commercial lithium ion batteries: A study
of 24 pouch cells.}
\textit{Journal of Power Sources}, 342, 589--597.

\bibitem{plett2015}
Plett, G.~L. (2015).
\textit{Battery Management Systems, Volume I: Battery Modeling.}
Artech House.

\bibitem{birkl2017}
Birkl, C.~R., Roberts, M.~R., McTurk, E., Bruce, P.~G., \& Howey, D.~A.
(2017).
\textit{Degradation diagnostics for lithium ion cells.}
\textit{Journal of Power Sources}, 341, 373--386.

\end{thebibliography}

\end{document}
